\documentclass[10pt]{article}
\usepackage[letterpaper,margin=.38in]{geometry}
\usepackage{amsmath,amssymb,multicol}
\pdfinfoomitdate=1\pdftrailerid{}
\pagestyle{empty}\setlength{\parindent}{0pt}\setlength{\parskip}{2pt}
\setlength{\columnsep}{13pt}\setlength{\jot}{0pt}
\setlength{\abovedisplayskip}{3pt}\setlength{\belowdisplayskip}{3pt}
\let\bigsum\sum\renewcommand{\sum}{\bigsum\nolimits}
\newcommand{\fl}[1]{\left\lfloor#1\right\rfloor}
\newcommand{\rg}[2]{\{#1,\ldots,#2\}}
\begin{document}
\begin{center}
{\large\bfseries A complete fifteen-variable consistency equation}\\[-1pt]
{\small Strictly positive integer unknowns; exact degree 157; a 66-operation evaluator.}
\end{center}\vspace{-8pt}
\begin{multicols}{2}
\fontsize{9}{9.5}\selectfont
The only unknowns are
\[
J,F,\alpha,z,f,h,i,T,s,w,\tau,\eta,\zeta,y,\gamma\in\mathbb Z_{>0}.
\]
All other symbols in the equation are polynomial abbreviations:
\[
\begin{gathered}
b=\beta J,\quad q=b+1,\quad C=q-F-\alpha,\\
X=wq,\quad Y=sq^3,\quad E=XY,\quad k=\eta+\zeta,\\
\ell=EkY,\quad c=kY+\eta,\quad a=E+Y,\\
H=4a+3,\quad\Delta=a^2+H,\quad D=X+ac+\gamma H,\\
R=[b(q-F)+\alpha](q^2-1)+(M+qN)J,\\
V=f^2(cT-R)-c,\quad S=i\Delta c^2,\\
P=\tau^2-\ell(\ell+k),\quad Q=D^2-\Delta c^2,\\
U=S^2(V^2-y^2)+y^2,\quad W=k-hE-R,\\
Z=(K+w)C+(q-F)-zb,\quad G=\Delta f^2-S^2.
\end{gathered}
\]
The single equation is
\[
\boxed{P\,Q\,U\,W\,Z\,G-\Delta=0.}
\]

\textbf{The four fixed integers.}
Every index, set, sum, product and floor below defines a constant, not
another unknown. Put
\[
\begin{gathered}
\rho=2^{47},\quad B=\rho^{2963},\quad\beta=B-1,\\
g(m,x)=\sum_{j=0}^{m-1}x^j,\quad
\Lambda=\rho^{1956}.
\end{gathered}
\]
\begin{align*}
M={}&\beta-g(951,\rho)\\
&-\rho^{960}g(26,\rho)-\rho^{1005},\\
N={}&\beta+4g(917,\rho)+(2^{44}-2)\Lambda\\
&+\rho^{960}g(9,\rho)+\rho^{987}g(8,\rho)
  +\rho^{1005},\\
K={}&\rho^{18}+\rho^{1005}+\rho^{2972}\\
&+\Lambda\Phi(1/\rho),
\end{align*}
where
\[
\begin{split}
\Phi(x)={}&\sum_{i=0}^{941}[1+2^{10}(1+w_i+2^8n_i\\
&+2^{17}p_i+2^{25}s_i)]x^i\\
&+2^{10}x^{969}(x^9+2^{17})g(8,2x)\\
&+2^{18}x^{942}(x^{18}+2^{17})g(9,2x).
\end{split}
\]
The rational intermediate in K cancels exactly; all four results are positive integers.

\textbf{The 942 indexed fixed tiles.}
Use the actions $u_j,v_j,d_j$ decoded below. For $d=0,1$, let
\[
\begin{gathered}
\theta(q,b)=2q+b+2,\quad
\mathcal S_d=\{u_j:d_j=d,\ 0\le j<430\},\\
r_{d,q}=205-199d+|\mathcal S_d\cap\rg0{q-1}|.
\end{gathered}
\]
An interval ending at $-1$ is empty. The table defines $t_i$ for
$0\le i\le941$; all index ranges are fixed:
\[
\setlength{\arraycolsep}{1.4pt}\begin{array}{c|l}
\text{index }i&\text{tile }t_i=(w_i,n_i,p_i,s_i)\\\hline
0&(2,435,3,1)\\
j+1\ (j=2q+b<430)&((1-d_j)r_{0,u_j},\theta(q,b),d_jr_{1,u_j},v_j)\\
433-2a+b\ (a,b=0,1)&(0,432a+b,0,432a+b)\\
\substack{423+2r_{d,q}+b\\q\in\mathcal S_d,\ d,b=0,1}&(dr_{d,q},b,(1-d)r_{d,q},\theta(q,b))\\
923+2a+b\ (a,b=0,1)&(0,432a+b,0,438)\\
927&(2,434,2,0)\\
928&(3,436,4,4)\\
929&(4,436,4,0)\\
496+t\ (t=434,435,436)&(1,438,1,t)\\
933+3a+b\ (a,b=0,1,2)&(a^2,437,b,437)
\end{array}
\]
\textbf{A closed decoder for the 430 actions.}
Put $\operatorname{digit}(A,b,k)=\fl{A/b^k}\bmod b$, where
$a\bmod b=a-b\fl{a/b}$. For $0\le q<215$, $b=0,1$, $j=2q+b$:
\[
\begin{gathered}
\mathsf{mask}=\mathcal A\bmod2^{636},\\
\mathsf{targets}=\fl{\mathcal A/2^{636}}\bmod2^{1675},\\
\mathsf{moves}=\fl{\mathcal A/2^{2311}},\\
\nu_j=\operatorname{digit}(\mathsf{mask},2,j),\quad \sigma_j=\sum_{k=0}^{j-1}\nu_k,\\
m_q=\operatorname{digit}(\mathsf{mask},2,430+q),\quad \delta_q=\sum_{k=0}^{q-1}m_k,\\
\omega_q=14+m_q(\operatorname{digit}(\mathsf{moves},16,\delta_q)-14),\\
\begin{aligned}
u_j={}&\nu_j(1+\sigma_j)+(1-\nu_j)\\
&\cdot\operatorname{digit}(\mathsf{targets},216,j-\sigma_j),
\end{aligned}\\
v_j=\operatorname{digit}(\omega_q,2,2b),\quad d_j=\operatorname{digit}(\omega_q,2,2b+1).
\end{gathered}
\]
Empty sums are zero. Here $(v_0,d_0,u_0)=(1,1,1)$; the constructor tile
uses this actual first target, including its state renaming.

\textbf{The integer $\mathcal A$ in base 10.}
Concatenate these thirteen lines, from top to bottom, to obtain its
781 decimal digits (there are no omitted digits):\par
\begingroup\fontsize{7.7}{8.2}\selectfont\ttfamily
5860725087898678188912673742632546324025172298353855272995683176\par
4969949406333919315323872375782477466725463011528567718294055076\par
3006124739855037353655204691391414186758883217097375044976429465\par
8270944142085505376842559029102693041085562540684690710604079067\par
3724763129630582855384181944414973901733570469173085418871671744\par
2048170942517613524282407203789879020689779835988780400299936946\par
5274600673267211886931165288781352441481875991871486564767549564\par
1498891563641180775282772787179257559378852323557013970506009282\par
7980034336115111037194604982010168556898406861821168844529953973\par
3329097089832109899184569780832475009429376795970989537956914018\par
6560751892408082293940602808596063442474573309002707517217784697\par
0890324537666784074933609456239656992860826945102893145422737549\par
3204418786549\par
\endgroup
\normalfont\fontsize{7.7}{8.2}\selectfont
The ordinary complemented-color, native and Pell arguments give
unsolvability iff $\mathrm{Con}(\mathrm{ZFC})$.
Native calibration and internal-ZFC formalization remain open in Lean.
This page is a complete finite specification, not an unsolvability proof.
\end{multicols}
\end{document}
